% vim: set foldmethod=marker foldlevel=0:

\documentclass[a4paper]{article}
\usepackage[UKenglish]{babel}

\usepackage{preamble}

\fancyhead[L]{MA2K4 Assignment 3}
\title{MA2K4 Numerical Methods and Computing, Assignment 3}
\colorlet{questionbodycolor}{green!50}

\begin{document}

\maketitle

\setlength{\parindent}{0em}
\setlength{\parskip}{1em}

% {{{ Q1
% \question{1}

% \begin{questionbody}
% Show that the Newton--Cotes quadrature rule of order $n > 1$ is exact for all polynomials of degree $n + 1$ if $n$ is even. % chktex 8

% % \textit{Hint}: For the integration interval $x \in [a, b]$, consider the polynomial \[
% % {\l( x - \df12 (a + b) \r)}^{n + 1}.
% % \]
% \end{questionbody}

% Suppose $n > 1$ is even, and consider the integral \[
% \intlim ab {{\l( x - \df12 (a + b) \r)}^{n + 1}} x.
% \]

% }}}

% {{{ Q2
\question{2}

\begin{questionbody}
Consider the Newton--Cotes quadrature rule of degree $n = 3$ on the interval $x \in [-1, 1]$, \[ % chktex 8
I_3(f) = \sum_{k=0}^3 \alpha_k f(x_k)
\] with $x_k = -1 + \df23 k$. Using the fact that this quadrature rule must be exact for all polynomials of degree 3, or otherwise, find the values of the four weights $\alpha_k$, for $k \in \{0, \dotsc, 3\}$.
\end{questionbody}

Let \[
f(x) = a x^3 + b x^2 + c x + d.
\] Then
\begin{align*}
I(f) &= \intlim {-1}1 {(a x^3 + b x^2 + c x + d)} x \\[0.5ex]
&= {\l[ \f a4 x^4 + \f b3 x^3 + \f c2 x^2 + dx \r]}_{-1}^1 \\[0.5ex]
&= \f a4 + \f b3 + \f c2 + d - \l( \f a4 - \f b3 + \f c2 - d \r) \\[0.5ex]
&= \f{2b}3 + 2d.
\end{align*}

The quadrature rule gives
\begin{align*}
I_3(f) &= \sum_{k=0}^3 \alpha_k f(x_k) \\[0.5ex]
&= \alpha_0 f(-1) + \alpha_1 f \l( -\f13 \r) + \alpha_2 f \l( \f13 \r) + \alpha_3 f(1) \\[0.5ex]
&= \alpha_0 (-a + b - c + d)
    + \alpha_1 \l( -\f1{27} a + \f19 b - \f13 c + d \r) \\
    &\qquad + \alpha_2 \l( \f1{27} a + \f19 b + \f13 c + d \r)
    + \alpha_3 (a + b + c + d) \\[0.5ex]
&= \l( -\alpha_0 - \f{\alpha_1}{27} + \f{\alpha_2}{27} + \alpha_3 \r) a
    + \l( \alpha_0 + \f{\alpha_1}9 + \f{\alpha_2}9 + \alpha_3 \r) b \\
    &\qquad + \l( -\alpha_0 - \f{\alpha_1}3 + \f{\alpha_2}3 + \alpha_3 \r) c
    + (\alpha_0 + \alpha_1 + \alpha_2 + \alpha_3) d.
\end{align*}

For these to be equivalent, we require
\begin{align*}
-\alpha_0 - \f{\alpha_1}{27} + \f{\alpha_2}{27} + \alpha_3 &= 0 \\
\alpha_0 + \f{\alpha_1}9 + \f{\alpha_2}9 + \alpha_3 &= \f23 \\
-\alpha_0 - \f{\alpha_1}3 + \f{\alpha_2}3 + \alpha_3 &= 0 \\
\alpha_0 + \alpha_1 + \alpha_2 + \alpha_3 &= 2
\end{align*}
We can rewrite this as a matrix equation and solve it with a computer. \[
\begin{pmatrix}
-1 & -\f1{27} & \f1{27} & 1 \\[0.5ex]
1 & \f19 & \f19 & 1 \\[0.5ex]
-1 & -\f13 & \f13 & 1 \\[0.5ex]
1 & 1 & 1 & 1
\end{pmatrix}
\begin{pmatrix}
\alpha_0 \\[0.5ex]
\alpha_1 \\[0.5ex]
\alpha_2 \\[0.5ex]
\alpha_3
\end{pmatrix}
= \begin{pmatrix}
0 \\[0.5ex]
\f23 \\[0.5ex]
0 \\[0.5ex]
2
\end{pmatrix}
\] This yields \[
\alpha_0 = \f14,\ \alpha_1 = \f34,\ \alpha_2 = \f34,\ \alpha_3 = \f14.
\]

% }}}

% {{{ Q3
\newquestion{3}

\begin{questionbody}
For the composite trapezoidal rule $I_{1, n}(f)$ and the composite Cavalieri--Simpson rule $I_{2, n}(f)$, show that \[ % chktex 8
I_{2, n}(f) = \f43 I_{1, 2n}(f) - \f13 I_{1, n}(f).
\]
\end{questionbody}

We know that
\begin{align*}
I_{1, n}(f) &= h \l( \f12 f(x_0) + f(x_1) + \dotsb + f(x_{n - 1}) + \f12 f(x_n) \r), \\[0.5ex]
I_{1, 2n}(f) &= h \l( \f12 f(x_0) + f(x_1) + \dotsb + f(x_{2n - 1}) + \f12 f(x_{2n}) \r), \\[0.5ex]
I_{2, n}(f) &= \f13 h \bigg( f(x_0) + 4 f(x_1) + 2 f(x_2) + 4 f(x_3) + \cdots \\
    &\qquad + 4 f(x_{2n - 3}) + 2 f(x_{2n - 2}) + 4 f(x_{2n - 1}) + f(x_{2n}) \bigg),
\end{align*}
but these definitions define $h$ differently and use different nodes. To normalise this, we instead use
\begin{align*}
I_{1, n}(f) &= h \l( \f12 f(x_0) + f(x_2) + \dotsb + f(x_{2n - 2}) + \f12 f(x_{2n}) \r), \\[0.5ex]
I_{1, 2n}(f) &= \f12 h \l( \f12 f(x_0) + f(x_1) + \dotsb + f(x_{2n - 1}) + \f12 f(x_{2n}) \r), \\[0.5ex]
I_{2, n}(f) &= \f16 h \bigg( f(x_0) + 4 f(x_1) + 2 f(x_2) + 4 f(x_3) + \cdots \\
    &\qquad + 4 f(x_{2n - 3}) + 2 f(x_{2n - 2}) + 4 f(x_{2n - 1}) + f(x_{2n}) \bigg),
\end{align*}
where $h = \f{b-a}{n}$.

Therefore
\begin{align*}
4 I_{1, 2n}(f) - I_{1, n}(f) &=
    2 h \l( \f12 f(x_0) + f(x_1) + \dotsb + f(x_{2n - 1}) + \f12 f(x_{2n}) \r) \\
    &\qquad - h \l( \f12 f(x_0) + f(x_2) + \dotsb + f(x_{2n - 2}) + \f12 f(x_{2n}) \r) \\[0.5ex]
&= h \bigg( \f12 f(x_0) + 2 f(x_1) + f(x_2) + 2 f(x_3) + \dotsb \\
    &\qquad + f(x_{2n - 2}) + 2 f(x_{2n - 1}) + \f12 f(x_{2n}) \bigg) \\[0.5ex]
&= \f12 h \bigg( f(x_0) + 4 f(x_1) + 2 f(x_2) + 4 f(x_3) + \cdots \\
    &\qquad + 4 f(x_{2n - 3}) + 2 f(x_{2n - 2}) + 4 f(x_{2n - 1}) + f(x_{2n}) \bigg) \\[0.5ex]
&= 3 I_{2, n}(f).
\end{align*}
Therefore $I_{2, n}(f) = \f43 I_{1, 2n}(f) - \f13 I_{1, n}(f)$ as required.

\hfill $\square$

% }}}

% {{{ Q4
\newquestion{4}

\begin{questionbody}
In this problem, show that it is not true that the Cavalieri--Simpson rule is necessarily more accurate that the trapezoidal rule, by constructing an explicit counterexample. % chktex 8

For this, calculate \[
\intlim 01 {\l( x^5 - \lambda x^4 \r)} x
\] and find the range of values of $\lambda \in \R$ for which the trapezoidal rule gives a more accurate estimate than the Cavalieri--Simpson rule. % chktex 8
\end{questionbody}

The exact value of the integral is
\begin{align*}
I(f) &= \intlim 01 {\l( x^5 - \lambda x^4 \r)} x \\[0.5ex]
&= {\l[ \f16 x^6 - \f\lambda5 x^5 \r]}_0^1 \\[0.5ex]
&= \f16 - \f\lambda5.
\end{align*}

The trapezoidal approximation and its error are
\begin{align*}
I_1(f) &= \f12 (1 - 0) (f(1) - f(0)) \\[0.5ex]
&= \f12 (1 - \lambda - 0) \\[0.5ex]
&= \f{1 - \lambda}2, \\[0.5ex]
E_1(f) &= \f16 - \f\lambda5 - \f{1 - \lambda}2 \\[0.5ex]
&= \f5{30} - \f{6 \lambda}{30} - \f{15 - 15 \lambda}{30} \\[0.5ex]
&= \f{-10 + 9 \lambda}{30}.
\end{align*}

The Cavalieri--Simpson approximation and its error are % chktex 8
\begin{align*}
I_2(f) &= \f{1 - 0}6 \l( f(0) + 4 f\l( \f12 (0 + 1) \r) + f(1) \r) \\[0.5ex]
&= \f16 \l( 0 + 4 \l( \f1{32} - \f\lambda{16} \r) + 1 - \lambda \r) \\[0.5ex]
&= \f16 \l( \f18 - \f\lambda4 + 1 - \lambda \r) \\[0.5ex]
&= \f{1 - 2 \lambda + 8 - 8 \lambda}{48} \\[0.5ex]
&= \f{9 - 10 \lambda}{48}, \\[0.5ex]
E_2(f) &= \f16 - \f\lambda5 - \f{9 - 10 \lambda}{48} \\[0.5ex]
&= \f{40}{240} - \f{48 \lambda}{240} - \f{45 - 50 \lambda}{240} \\[0.5ex]
&= \f{-5 + 2 \lambda}{240}.
\end{align*}

We want $|E_1(f)| < |E_2(f)|$, so we will first find where $E_1(f) = 0$ so we can centre our search.
\begin{align*}
\f{-10 + 9 \lambda}{30} &= 0 \\
9 \lambda &= 10 \\
\lambda &= \f{10}9
\end{align*}
so we know that our final range must include $\f{10}9$.

Each of $E_1(f)$ and $E_2(f)$ may be positive or negative, so we have 4 inequalities.
\begin{align*}
E_1(f) &< E_2(f) \\
\f{-10 + 9 \lambda}{30} &< \f{-5 + 2 \lambda}{240} \\
24(-10 + 9 \lambda) &< 3(-5 + 2 \lambda) \\
-240 + 216 \lambda &< -15 + 6 \lambda \\
210 \lambda &< 225 \\
\lambda &< \f{15}{14}
\end{align*}

\begin{align*}
E_1(f) &< -E_2(f) \\
\f{-10 + 9 \lambda}{30} &< \f{5 - 2 \lambda}{240} \\
-240 + 216 \lambda &< 15 - 6 \lambda \\
222 \lambda &< 255 \\
\lambda &< \f{84}{74}
\end{align*}

\begin{align*}
-E_1(f) &< E_2(f) \\
\f{10 - 9 \lambda}{30} &< \f{-5 + 2 \lambda}{240} \\
240 - 216 \lambda &< -15 + 6 \lambda \\
255 &< 222 \lambda \\
\f{84}{74} &< \lambda
\end{align*}

\begin{align*}
-E_1(f) &< -E_2(f) \\
\f{10 - 9 \lambda}{30} &< \f{5 - 2 \lambda}{240} \\
240 - 216 \lambda &< 15 - 6 \lambda \\
225 &< 210 \lambda \\
\f{15}{14} &< \lambda
\end{align*}

Since we need $\f{10}9$ to be in the range, we conclude that the desired range is \[
\f{15}{14} < \lambda < \f{85}{74}.
\] Testing shows that the trapezoidal rule is in fact better that Cavalieri--Simpson for this integral with $\lambda$ in that range. % chktex 8

% }}}

% {{{ Q5
\newquestion{5}

\begin{questionbody}
Derive the 3-node Gauss quadrature on the interval $[-1, 1]$ by finding the three weights $\alpha_i$ and nodes $x_i$ for $i \in \{0, 1, 2\}$. Then, use the rescaling formula \[
G_n(f) = \f12 (b - a) \sum_{i=0}^n \alpha_i f\l( \f12 (a + b) + x_i \f12 (b - a) \r)
\] to write down the 3-node Gauss quadrature on an arbitrary interval $[a, b]$.
\end{questionbody}

Our maximal degree of precision is $2n + 1 = 5$, so we need
\begin{align*}
G_2(1) &= 2 \\
G_2(x) &= 0 \\
G_2(x^2) &= \f23 \\
G_2(x^3) &= 0 \\
G_2(x^4) &= \f25 \\
G_2(x^5) &= 0.
\end{align*}
At the same time, \[
G_2(f) = \alpha_0 f(x_0) + \alpha_1 f(x_1) + \alpha_2 f(x_2).
\]

This gives
\begin{align*}
\alpha_0 + \alpha_1 + \alpha_2 &= 2 \\
\alpha_0 x_0 + \alpha_1 x_1 + \alpha_2 x_2 &= 0 \\
\alpha_0 x_0^2 + \alpha_1 x_1^2 + \alpha_2 x_2^2 &= \f23 \\
\alpha_0 x_0^3 + \alpha_1 x_1^3 + \alpha_2 x_2^3 &= 0 \\
\alpha_0 x_0^4 + \alpha_1 x_1^4 + \alpha_2 x_2^4 &= \f25 \\
\alpha_0 x_0^5 + \alpha_1 x_1^5 + \alpha_2 x_2^5 &= 0.
\end{align*}
By symmetry, we require $\alpha_0 = \alpha_2$, $x_0 = -x_2$, and $x_1 = 0$, so we get
\begin{align*}
2 \alpha_0 + \alpha_1 &= 2 \\
0 &= 0 \\
2 \alpha_0 x_0^2 &= \f23 \\
0 &= 0 \\
2 \alpha_0 x_0^4 &= \f25 \\
0 &= 0
\end{align*}
since for odd $k$, $\alpha_0 x_0^k + \alpha_2 x_2^k = 0$.

Then we compute the solution and find
\begin{align*}
x_0^2 &= \f35 \\[0.5ex]
2 \alpha_0 &= \f23 \cdot \f53 \\[0.5ex]
&= \f{10}9 \\[0.5ex]
\alpha_0 &= \f59 \\[0.5ex]
\alpha_1 &= 2 - \f{10}9 \\[0.5ex]
&= \f89
\end{align*}
So we get $x_0 = -\sqrt{\df35}$, $x_1 = 0$, $x_2 = \sqrt{\df35}$, $\alpha_0 = \df59$, $\alpha_1 = \df89$, and $\alpha_2 = \df59$.
% This is correct.

Therefore
\begin{align*}
G_2(f) &= \f12 (b - a) \sum_{i=0}^2 \alpha_i f\l( \f12 (a + b) + x_i \f12 (b - a) \r) \\[0.5ex]
&= \f12 (b - a) \bigg(
    \alpha_0 f\l( \f12 (a + b) + \f12 x_0 (b - a) \r) \\
    &\qquad + \alpha_1 f\l( \f12 (a + b) + \f12 x_1 (b - a) \r) \\
    &\qquad + \alpha_2 f\l( \f12 (a + b) + \f12 x_2 (b - a) \r)
\bigg) \\[0.5ex]
&= \f12 (b - a) \Bigg(
    \f59 f\l( \f12 (a + b) - \f12 \sqrt{\f35} (b - a) \r) \\
    &\qquad + \f89 f\l( \f12 (a + b) \r)
    + \f59 f\l( \f12 (a + b) + \f12 \sqrt{\f35} (b - a) \r)
\Bigg) \\[0.5ex]
&= \f1{18} (b - a) \Bigg(
    5 f\l( \f12 (a + b) - \f{\sqrt 3}{2 \sqrt 5} (b - a) \r) \\
    &\qquad + 8 f\l( \f12 (a + b) \r)
    + 5 f\l( \f12 (a + b) + \f{\sqrt 3}{2 \sqrt 5} (b - a) \r)
\Bigg).
\end{align*}

% }}}

% {{{ Q6
\newquestion{6}

\begin{questionbody}
Consider the initial value problem \[
y'(t) = {(y(t))}^{1/5}, \quad y(0) = 0.
\]
\begin{enumerate}[(a)]
\item Show that \[
y(t) = {\l( \f45 t \r)}^{5/4}
\] is a solution of the above IVP\@.

\item Show that the forward Euler method $u_{n+1} = u_n + h u_n^{1/5}$ with $u_0 = y(0)$ fails to approximate this solution. Give a reason for why this is the case.
\end{enumerate}
\end{questionbody}

\subsection{~} % 6.a

Clearly $y(0) = 0^{5/4} = 0$ as required, so we just need to check the derivative.
\begin{align*}
y'(t) &= \f54 \f45 {\l( \f45 t \r)}^{1/4} \\
&= {\l( \f45 t \r)}^{1/4} \\
&= {\l( {\l( \f45 t \r)}^{5/4} \r)}^{1/5} \\
&= {(y(t))}^{1/5}
\end{align*}
as required, so this is a solution to the IVP\@.

\subsection{~} % 6.b

The forward Euler method gives
\begin{align*}
u_0 &= 0 \\
u_1 &= 0 + h \cdot 0^{1/5} \\
&= 0 \\
u_2 &= 0 + h \cdot 0^{1/5} \\
&= 0 \\
&\ \ \vdots
\end{align*}
whereas of course $y(t) = {\l(\f45 t\r)}^{5/4}$ is not always 0.

Since this IVP is autonomous and starts at 0, each step of the Euler method stays at 0, so the process can never `escape'.

% }}}

\end{document}
